% Chapter 7 -- Conclusion and Future Work
% Status: basic draft.
\chapter{Conclusion and Future Work}
\label{ch:conclusion}

\section{Conclusion}
\label{sec:conc-conclusion}

This thesis asked where personalisation should live in a cross-device federated recommender,
what each placement costs in privacy, and how stateful clients must be evaluated. Under one
locked protocol on MovieLens 1M (\NClients{} clients, $C=\ParticipationC$, \NRounds{} rounds,
shared candidate lists), the placement of the personal part dominated every other design
choice. Keeping the user vector on the device improved \ndcg{} by
\GainPersonalizedVsBaselinePct\% over averaging it with nothing else changed; a private scoring
head on a shared item table reached \NdcgPfedrec, and on that model the private head was worth
\DeltaHeadGlobalHead{} against the population head, with no blend between the two recovering
the loss. The thesis's own adaptive method, which blends the user vector with a population
prototype and adds a per-user neural scorer, was reported as a negative result with three
measured or identified causes: the blend pulls taste toward the mean where the prototype is
least adequate, the private network is oversized for the data, and the upload of the user vector
opens a leakage channel that the other placements do not have.

The comparison with the published PFedRec number was treated as a measurement problem. A
pre-registered re-score showed that client-side post-tuning of the item table lowers \ndcg{}
under honest negative sampling and raises it only when the held-out item is excluded from the
negatives as in the released code, a difference of \PtContrastEOne{} carried entirely by the
\PtContrastEOneNonZero{} users whose test item was drawn. The thesis reports its number with
the shared item table and states the measured cause in its limitations. Finally, the work
established the safeguards without which none of these numbers could be trusted: a
cold-evaluation error that reduced \ndcg{} \ColdEvalRatio-fold was found and fixed, cache-miss
counters and validity sidecars were added to every run, and an offline re-scoring tool that
reproduces the federated evaluation exactly became the basis of every analysis.
\pending{one sentence on the privacy measurements and the stale-head cost once they are in}.

\section{Future Work}
\label{sec:conc-future}

\begin{itemize}[leftmargin=2em]
  \item \textbf{Private item table with row-wise blending.} A client-held item table
        $\Qu$ that is blended with the global table only on the rows of items the user
        interacted with was proposed during this work and withdrawn after analysis: the private
        part would then live only on items the user has already seen and could not reach the
        items to be recommended. A formulation that propagates the private adjustment to
        unseen items (for example through item--item similarity) remains open, and its
        evaluation must use the clean negative pool of Section~\ref{sec:method-posttune}.
  \item \textbf{Secure aggregation and a measured privacy budget.} The item-table update
        pattern is the leakage channel common to all placements; implementing secure
        aggregation \cite{bonawitz2017secagg} in the framework and measuring what the server
        can still infer would close the gap between the threat model and the deployment.
  \item \textbf{Seeds and a validation item.} Two further seeds for the three close cells and a
        protocol with a validation item would turn the single-seed point estimates into
        intervals and remove the best-round inflation.
  \item \textbf{Participation fraction.} A PFedRec run at $C=1$ under the thesis protocol
        would separate the share of the reproduction gap due to participation from the share
        due to the deviations K1--K3.
  \item \textbf{A cold-start population.} The adaptive design should be re-evaluated where its
        assumptions hold, on a dataset with genuinely cold users or with a train/test shift per
        user, before it is abandoned.
\end{itemize}
